\documentclass[twocolumn]{article}

% ---- Packages ----
\usepackage[utf8]{inputenc}
\usepackage[T1]{fontenc}
\usepackage{lmodern}
\usepackage{microtype}
\usepackage{graphicx}
\usepackage{booktabs}
\usepackage{hyperref}
\usepackage{xcolor}
\usepackage[margin=2cm]{geometry}
\usepackage{enumitem}
\usepackage{amsmath}

\hypersetup{
  colorlinks=true,
  linkcolor=blue!60!black,
  citecolor=blue!60!black,
  urlcolor=blue!60!black
}

% ---- Title ----
\title{BIME: A Browser-Based Molecular Editor with Built-in Substructure Intelligence}

\author{
  Syed Asad Rahman\thanks{Corresponding author: \texttt{s9asad@gmail.com}} \\
  BioInception PVT LTD, Cambridge, UK
}

\date{}

\begin{document}

\maketitle

% ====================================================================
\begin{abstract}
BIME (BioInception Molecular Editor) is a browser-based molecular
structure editor written in pure JavaScript with zero runtime
dependencies.  Its three principal contributions are: (1)~a built-in
maximum common substructure (MCS) engine and molecular fingerprint
toolkit that run entirely in the browser, eliminating server
round-trips; (2)~an implementation of IUPAC 2013 Cahn--Ingold--Prelog
(CIP) Rules~1--4a for R/S and E/Z stereodescriptor assignment using
digraph expansion with phantom nodes; and (3)~a zero-dependency,
single-file architecture that enables fully offline deployment for
education and resource-constrained settings.  BIME embeds the SMSD Pro
engine (v6.6.1) for MCS detection via McSplit partition refinement,
substructure search via VF2++, and six fingerprint similarity metrics.
A layout engine with 45 scaffold templates, crossing reduction, and
SMACOF stress majorisation produces publication-quality 2D coordinates.
BIME is freely available under the Apache License~2.0 at
\url{https://asad.github.io/bime-dist/}.
\end{abstract}

\noindent\textbf{Keywords:} molecular editor, maximum common
substructure, substructure search, molecular fingerprints,
Cahn--Ingold--Prelog, SMILES, JavaScript, SVG

% ====================================================================
\section{Background}

Browser-based molecular editors are essential components of modern
cheminformatics web applications, enabling researchers to draw, edit,
and query chemical structures without installing desktop
software~\citep{Bienfait2013,Ertl2010}.  Several capable editors
exist---JSME~\citep{Bienfait2013}, Ketcher~\citep{Ketcher2024}, and
ChemDoodle Web Components~\citep{Burger2015}---yet they share a common
architectural pattern: drawing and cheminformatics intelligence are
decoupled.  MCS computation, fingerprint generation, and stereochemistry
assignment typically require a server-side toolkit such as
RDKit~\citep{Landrum2006} or CDK~\citep{Steinbeck2003}, introducing
latency and precluding fully offline workflows.

\paragraph{WebAssembly approaches.}
Recent efforts have brought server-side toolkits to the browser via
WebAssembly (WASM).  RDKit.js~\citep{RDKitJS2024} compiles the C++
RDKit core to WASM, providing substructure search, fingerprints, and
coordinate generation in the client.  OpenChemLib~JS similarly
cross-compiles a Java toolkit.  These approaches offer mature algorithm
coverage but carry trade-offs: WASM modules add 5--15\,MB to the
initial download, require asynchronous loading, and cannot run from
\texttt{file://} URIs without a local server due to CORS restrictions.
For applications requiring a small footprint, instant startup, and
true offline use (e.g., loading an HTML file from a USB drive in a
field laboratory), a pure-JavaScript solution avoids these constraints.

This gap is particularly acute in educational settings, field
laboratories, and resource-constrained networks where reliable server
access cannot be assumed.  Furthermore, many web applications need only
lightweight substructure queries or pairwise similarity scores and
cannot justify the infrastructure cost of a full cheminformatics server.

BIME addresses this gap by embedding a complete substructure search and
similarity engine---SMSD Pro~v6.6.1---directly into the editor.  SMSD
Pro is an independent JavaScript implementation that builds on the
algorithmic ideas first described in the original Small Molecule
Subgraph Detector~\citep{Rahman2009}, substantially extended with
modern algorithms including McSplit partition
refinement~\citep{McCreesh2017}, VF2++ backtracking~\citep{Juttner2018},
and multi-level coverage-driven MCS funnelling.  The result is a single
JavaScript file that delivers drawing, I/O, stereochemistry, layout,
MCS, fingerprints, and similarity search with no runtime dependencies.

% ====================================================================
\section{Implementation}

\subsection{Architecture}

BIME is implemented in approximately 25\,000 lines of ES5 JavaScript,
organised into 23 self-contained modules following the Immediately
Invoked Function Expression (IIFE) pattern.  Each module exposes its
public API on the \texttt{window} object.  The modules are concatenated
into a single distributable file; no transpiler, bundler, or package
manager is required.  This design ensures compatibility with all modern
browsers (Chrome, Firefox, Safari, Edge) and allows BIME to run from
\texttt{file://} URIs with no web server.

\subsection{Molecular Graph Model}

The core data structure (\texttt{Molecule.js}) represents molecules as
an atom--bond graph with 2D coordinates.  Atoms carry element symbol,
charge, isotope, chirality (\texttt{@}/\texttt{@@}), hydrogen count, and
atom-map number.  Bonds carry type (single, double, triple, aromatic),
stereo (wedge, dash), and directional markers (\texttt{/},
\texttt{\char`\\}) for E/Z specification following the
SMILES~\citep{Weininger1988} and OpenSMILES~\citep{OpenSMILES2024}
conventions.  Reaction SMILES are supported
natively via \texttt{>>} separation of reactant and product components.

\subsection{SMSD Pro Engine}

Six modules implement the SMSD Pro engine:

\begin{description}[nosep,leftmargin=1em]
\item[\texttt{SmsdGraph}] Converts the editor's molecular graph into
  flat adjacency arrays with atom labels, bond labels, degree sequences,
  and Neighbourhood Label Frequency (NLF) vectors for fast
  compatibility checking.
\item[\texttt{SmsdVF2}] VF2++~\citep{Juttner2018} substructure
  isomorphism with FASTiso ordering for small molecules ($<$30 atoms)
  and VF3-Light ordering for larger targets, NLF-1/2 pruning, and a
  greedy $O(N)$ fast-path.  The algorithm extends the original
  VF2~\citep{Cordella2004} with improved candidate selection.
\item[\texttt{SmsdMCS}] A coverage-driven MCS funnel with seven+
  algorithmic levels: label-frequency upper bound, linear-chain and tree
  fast-paths, greedy probe, containment check, augmenting-path
  refinement, McSplit~\citep{McCreesh2017} partition refinement with
  RRSplit maximality, Bron--Kerbosch~\citep{Bron1973} with
  Tomita pivoting~\citep{Tomita2006} and $k$-core reduction, and
  McGregor bond-grow extension.
\item[\texttt{SmsdRings}] Smallest Set of Smallest Rings (SSSR)
  via Horton candidate generation and GF(2) Gaussian elimination,
  Relevant Cycle Basis, and Unique Ring Families.
\item[\texttt{SmsdBatch}] Batch substructure search, batch MCS, RASCAL
  fingerprint pre-screening, path-based fingerprints,
  ECFP/FCFP~\citep{Rogers2010} circular fingerprints following the
  Morgan~\citep{Morgan1965} extended-connectivity algorithm,
  topological torsion fingerprints, and six
  similarity metrics (Tanimoto, Dice, cosine, Soergel, count-Tanimoto,
  count-Dice).
\item[\texttt{SmsdLayout}] SSSR-aware layout utilities including ring
  system ordering, simulated-annealing crossing reduction, PrEd-inspired
  force-directed refinement~\citep{Bertault2000}, and SMACOF stress
  majorisation~\citep{deLeeuw1977}.
\end{description}

\subsection{CIP Stereochemistry}

The \texttt{CipStereo} module implements Rules~1--4a of the IUPAC 2013
CIP recommendations~\citep{Favre2013}, building on the classical
framework of Prelog and Helmchen~\citep{Prelog1982}.  A digraph is
constructed by BFS expansion from each stereocentre with phantom-node
duplication for double/triple bonds and ring back-edges.  Priority
comparison proceeds level by level.  Rule~3 resolves ties via
seqCis/seqTrans descriptors ($Z > E > $ nonstereogenic); Rule~4a applies
a two-pass like/unlike comparison for chiral centres.  Depth and node
limits (\texttt{MAX\_DEPTH}~=~30, \texttt{MAX\_NODES}~=~5000) guarantee
termination on pathological inputs.  R/S labels are assigned to
tetrahedral centres and E/Z labels to restricted double bonds.

\subsection{2D Layout}

The layout engine (\texttt{Layout.js}) generates publication-quality
coordinates following the Structure Diagram Generation (SDG)
architecture~\citep{Helson1999}: (1)~SSSR ring perception,
(2)~ring-system fusion with regular-polygon placement, edge-reflection
for fused rings, and spiro-junction handling, (3)~120-degree zigzag
chain extension, (4)~substituent fan placement, and (5)~iterative
overlap resolution using a spatial hash grid for $O(N)$ collision
detection.  Forty-five pre-computed scaffold
templates~\citep{Bemis1996} provide instant coordinates for common ring
systems (see Figure~\ref{fig:workbench}).

\subsection{SMARTS Support}

A recursive-descent SMARTS parser (\texttt{SmartsParser.js}) produces
query molecules carrying per-atom constraint arrays (wildcard, atomic
number, ring membership, degree, valence, hydrogen count, charge,
aromaticity, logical AND/OR/NOT, recursive SMARTS).  The VF2-based
matcher (\texttt{SmartsMatch.js}) evaluates these constraints during
backtracking.  Round-trip fidelity is maintained via
\texttt{SmartsWriter.js}.

% ====================================================================
\section{Features}

BIME provides a full drawing toolkit: single, double, and triple bonds;
wedge/dash stereo bonds; chain drawing with angle snapping at 30-degree
intervals; ring templates (3--8-membered rings, fused bicyclics); atom
label editing with auto-valence hydrogen adjustment; charge and isotope
modification; and an unlimited undo/redo stack (100 levels).

Import and export are supported for SMILES (OpenSMILES
specification~\citep{OpenSMILES2024} with full stereochemistry),
reaction SMILES (\texttt{>>}), SMARTS, MOL V2000, and MOL V3000
(automatic fallback for molecules exceeding 999 atoms).  The canonical
SMILES writer uses Morgan-algorithm~\citep{Morgan1965} ordering and
H\"uckel $4n+2$ aromatic perception.  Publication-quality SVG and PNG
export with configurable dimensions, padding, background colour, and
font are supported.

A library of 552 structures---covering common drugs, natural products,
amino acids, nucleotides, vitamins, and reagents---sourced from PubChem
and ChEMBL provides instant access to frequently used molecules.
Structures were obtained as SMILES strings, which represent
non-copyrightable molecular facts.

% ---- Figures ----

\begin{figure}[t]
\centering
\fbox{\parbox{0.9\columnwidth}{\centering\small
  [Screenshot of BIME workbench showing the editor canvas,\\
  toolbar, and SMSD Pro panel with MCS/fingerprint controls]\\[4pt]
  Interactive demo:
  \url{https://asad.github.io/bime-dist/workbench.html}}}
\caption{The BIME workbench interface.  The left panel provides drawing
  tools and ring templates; the right panel exposes SMSD Pro
  functionality including MCS computation, substructure search, and
  fingerprint similarity scoring.}
\label{fig:workbench}
\end{figure}

\begin{figure}[t]
\centering
\fbox{\parbox{0.9\columnwidth}{\centering\small
  [MCS result: two molecules displayed side by side with the\\
  common substructure highlighted in colour on both structures]}}
\caption{Maximum common substructure (MCS) result.  The common
  substructure atoms and bonds are highlighted on both the query and
  target molecules.  The MCS was computed entirely in the browser
  using the McSplit/Bron--Kerbosch funnel.}
\label{fig:mcs}
\end{figure}

\begin{figure}[t]
\centering
\fbox{\parbox{0.9\columnwidth}{\centering\small
  [Morphine layout: left --- before crossing reduction (10 crossings);\\
  right --- after two-phase optimisation (2 crossings)]}}
\caption{Crossing reduction on morphine (5~rings, 40~atoms).
  Left: initial layout with 10~bond crossings.
  Right: after two-phase simulated-annealing optimisation,
  crossings are reduced to~2.}
\label{fig:crossings}
\end{figure}

% ====================================================================
\section{Results and Discussion}

\subsection{Performance}

MCS computation on 200 random ChEMBL drug pairs (20--50 heavy atoms)
in Chrome~124 on a 2.6\,GHz Apple M2 MacBook Pro yielded a median time
of 38\,ms and a 95th-percentile time of 210\,ms.  The multi-level
funnel architecture ensures that simple cases (substructure containment,
identical scaffolds) are resolved at early levels without engaging the
more expensive McSplit or Bron--Kerbosch algorithms.  The adaptive
timeout (default 10\,s) guarantees responsiveness even for worst-case
inputs.

Full 2D layout generation for taxol (113 atoms, 11 rings) completes in
66\,ms, including SSSR perception, ring-system fusion, chain layout,
and overlap resolution.  Template matching bypasses the full layout
pipeline for the 45 most common scaffolds, providing sub-millisecond
coordinate generation.  Systematic benchmarking against standard
datasets (e.g., the GDB and COCONUT collections) is planned for a
future study.

\subsection{Crossing Reduction}

The simulated-annealing crossing reduction stage, combined with
PrEd-inspired force-directed refinement, substantially improves diagram
clarity (Figure~\ref{fig:crossings}).  For morphine (5 rings, 40 atoms),
bond crossings are reduced from 10 to 2.  A two-phase strategy first
optimises ring-system orientations globally, then refines individual
ring flips, avoiding local minima that trap single-phase approaches.

\subsection{CIP Correctness}

The CIP implementation was validated against the canonical examples in
the IUPAC 2013 recommendations~\citep{Favre2013} and a curated set of
challenging molecules including allenes, atropisomers (via E/Z on
restricted bonds), and spiro centres.  The digraph-based approach with
phantom-node expansion correctly handles cases where path-based
heuristics used by simpler implementations fail, such as bridged
bicyclic systems with multiple stereocentres.

\subsection{Fingerprint Quality}

ECFP and FCFP circular fingerprints~\citep{Rogers2010} follow the
Morgan/extended connectivity algorithm~\citep{Morgan1965} with
configurable radius.  Path-based fingerprints enumerate all paths up to
a specified length and hash atom and bond labels.  Preliminary
comparison of Tanimoto similarity scores computed by BIME against
RDKit-computed values on a set of kinase inhibitors showed strong
agreement; formal validation on standard benchmark sets is underway.

% ====================================================================
\section{Limitations}
\label{sec:limitations}

Several limitations should be noted.  First, the CIP implementation
covers Rules~1--4a; Rules~4b--5 (pseudoasymmetric centres and the
full reflection/rotation invariance criteria) are not yet implemented
and will be addressed in a future release.  Second, the crossing
reduction heuristic does not fully eliminate crossings for all bridged
polycyclic systems---morphine, for example, retains 2~bond crossings
after optimisation.  Third, BIME does not use WebAssembly acceleration;
for molecules exceeding approximately 200 heavy atoms, server-side
toolkits such as RDKit~\citep{Landrum2006} or
CDK~\citep{Steinbeck2003} will offer superior performance.  Fourth, the
distributed JavaScript is obfuscated for intellectual-property
protection; readable source code is available from the author on
request for academic verification purposes.

% ====================================================================
\section{Conclusions}

BIME is a lightweight molecular editor with zero runtime dependencies
that brings substructure intelligence---MCS, fingerprints, similarity
search, and CIP stereochemistry---directly to the browser.  Its
single-file deployment model makes it suited to offline, educational,
and resource-constrained settings.  Future work will extend CIP coverage
to Rules~4b--5, add pharmacophore perception, and integrate WebAssembly
acceleration for molecules exceeding 200 heavy atoms.

% ====================================================================
\section*{Availability and Requirements}

\begin{description}[nosep,leftmargin=1em,font=\normalfont\bfseries]
\item[Project name:] BIME --- BioInception Molecular Editor
\item[Live demo:] \url{https://asad.github.io/bime-dist/workbench.html}
\item[Distribution:] \url{https://github.com/asad/bime-dist}
\item[Operating system(s):] Platform independent (browser-based)
\item[Programming language:] JavaScript (ES5)
\item[Other requirements:] Any modern web browser (Chrome, Firefox, Safari, Edge)
\item[License:] Apache License 2.0 (freely available)
\item[Source code:] Readable source available from the author on request for academic verification
\item[Any restrictions to use by non-academics:] None
\end{description}

% ====================================================================
\section*{Abbreviations}

\noindent
BIME: BioInception Molecular Editor;
CIP: Cahn--Ingold--Prelog;
ECFP: Extended-Connectivity Fingerprint;
FCFP: Functional-Class Fingerprint;
IIFE: Immediately Invoked Function Expression;
MCS: Maximum Common Substructure;
NLF: Neighbourhood Label Frequency;
SDG: Structure Diagram Generation;
SMARTS: SMILES Arbitrary Target Specification;
SMILES: Simplified Molecular-Input Line-Entry System;
SMSD: Small Molecule Subgraph Detector;
SSSR: Smallest Set of Smallest Rings;
SVG: Scalable Vector Graphics;
VF2++: an improved subgraph isomorphism algorithm;
WASM: WebAssembly

% ====================================================================
\section*{Acknowledgements}

The author acknowledges the algorithmic foundations laid by the original
Small Molecule Subgraph Detector~\citep{Rahman2009} and the SMSD
community (2008--2018).  SMSD Pro is an independent implementation and
substantial extension by BioInception PVT LTD, not a fork or derivative
of the original Java codebase.  The author also acknowledges the
broader cheminformatics open-source ecosystem, including the CDK and
RDKit projects, whose published algorithms have informed the field.

% ====================================================================
\begin{thebibliography}{99}

\bibitem[{Rahman et~al.(2009)}]{Rahman2009}
Rahman, S.A., Bashton, M., Holliday, G.L., Schrader, R., Thornton, J.M.
\newblock Small Molecule Subgraph Detector (SMSD) toolkit.
\newblock \emph{Journal of Cheminformatics}, 1:12, 2009.
\newblock \href{https://doi.org/10.1186/1758-2946-1-12}{doi:10.1186/1758-2946-1-12}

\bibitem[{Bienfait and Ertl(2013)}]{Bienfait2013}
Bienfait, B., Ertl, P.
\newblock JSME: a free molecule editor in JavaScript.
\newblock \emph{Journal of Cheminformatics}, 5:24, 2013.
\newblock \href{https://doi.org/10.1186/1758-2946-5-24}{doi:10.1186/1758-2946-5-24}

\bibitem[{Ertl(2010)}]{Ertl2010}
Ertl, P.
\newblock Molecular structure input on the web.
\newblock \emph{Journal of Cheminformatics}, 2:1, 2010.
\newblock \href{https://doi.org/10.1186/1758-2946-2-1}{doi:10.1186/1758-2946-2-1}

\bibitem[{Burger et~al.(2015)}]{Burger2015}
Burger, M.C.
\newblock ChemDoodle Web Components: HTML5 toolkit for chemical graphics,
  interfaces, and informatics.
\newblock \emph{Journal of Cheminformatics}, 7:35, 2015.
\newblock \href{https://doi.org/10.1186/s13321-015-0085-3}{doi:10.1186/s13321-015-0085-3}

\bibitem[{EPAM(2024)}]{Ketcher2024}
EPAM Systems.
\newblock Ketcher: web-based molecular editor.
\newblock \url{https://github.com/epam/ketcher}, 2024.

\bibitem[{Bertault(2000)}]{Bertault2000}
Bertault, F.
\newblock A force-directed algorithm that preserves edge-crossing properties.
\newblock \emph{Information Processing Letters}, 74(1--2):7--13, 2000.
\newblock \href{https://doi.org/10.1016/S0020-0190(00)00028-8}{doi:10.1016/S0020-0190(00)00028-8}

\bibitem[{de~Leeuw(1977)}]{deLeeuw1977}
de~Leeuw, J.
\newblock Applications of convex analysis to multidimensional scaling.
\newblock In \emph{Recent Developments in Statistics}, pages 133--145.
  North-Holland, 1977.

\bibitem[{Favre and Powell(2013)}]{Favre2013}
Favre, H.A., Powell, W.H.
\newblock \emph{Nomenclature of Organic Chemistry: IUPAC Recommendations and
  Preferred Names 2013}.
\newblock Royal Society of Chemistry, 2013.
\newblock \href{https://doi.org/10.1039/9781849733069}{doi:10.1039/9781849733069}

\bibitem[{Helson(1999)}]{Helson1999}
Helson, H.E.
\newblock Structure diagram generation.
\newblock \emph{Reviews in Computational Chemistry}, 13:313--398, 1999.
\newblock \href{https://doi.org/10.1002/9780470125908.ch6}{doi:10.1002/9780470125908.ch6}

\bibitem[{McCreesh and Prosser(2017)}]{McCreesh2017}
McCreesh, C., Prosser, P.
\newblock A partitioning algorithm for maximum common subgraph problems.
\newblock In \emph{Proceedings of the 26th International Joint Conference on
  Artificial Intelligence (IJCAI)}, pages 712--719, 2017.
\newblock \href{https://doi.org/10.24963/ijcai.2017/99}{doi:10.24963/ijcai.2017/99}

\bibitem[{J\"uttner and Madarasi(2018)}]{Juttner2018}
J\"uttner, A., Madarasi, P.
\newblock {VF2++}---An improved subgraph isomorphism algorithm.
\newblock \emph{Discrete Applied Mathematics}, 242:69--81, 2018.
\newblock \href{https://doi.org/10.1016/j.dam.2018.02.018}{doi:10.1016/j.dam.2018.02.018}

\bibitem[{Tomita et~al.(2006)}]{Tomita2006}
Tomita, E., Tanaka, A., Takahashi, H.
\newblock The worst-case time complexity for generating all maximal cliques and
  computational experiments.
\newblock \emph{Theoretical Computer Science}, 363(1):28--42, 2006.
\newblock \href{https://doi.org/10.1016/j.tcs.2006.06.015}{doi:10.1016/j.tcs.2006.06.015}

\bibitem[{Morgan(1965)}]{Morgan1965}
Morgan, H.L.
\newblock The generation of a unique machine description for chemical
  structures---a technique developed at Chemical Abstracts Service.
\newblock \emph{Journal of Chemical Documentation}, 5(2):107--113, 1965.
\newblock \href{https://doi.org/10.1021/c160017a018}{doi:10.1021/c160017a018}

\bibitem[{Rogers and Hahn(2010)}]{Rogers2010}
Rogers, D., Hahn, M.
\newblock Extended-connectivity fingerprints.
\newblock \emph{Journal of Chemical Information and Modeling}, 50(5):742--754, 2010.
\newblock \href{https://doi.org/10.1021/ci100050t}{doi:10.1021/ci100050t}

\bibitem[{Weininger(1988)}]{Weininger1988}
Weininger, D.
\newblock {SMILES}, a chemical language and information system. 1.~Introduction
  to methodology and encoding rules.
\newblock \emph{Journal of Chemical Information and Computer Sciences},
  28(1):31--36, 1988.
\newblock \href{https://doi.org/10.1021/ci00057a005}{doi:10.1021/ci00057a005}

\bibitem[{Cordella et~al.(2004)}]{Cordella2004}
Cordella, L.P., Foggia, P., Sansone, C., Vento, M.
\newblock A (sub)graph isomorphism algorithm for matching large graphs.
\newblock \emph{IEEE Transactions on Pattern Analysis and Machine Intelligence},
  26(10):1367--1372, 2004.
\newblock \href{https://doi.org/10.1109/TPAMI.2004.75}{doi:10.1109/TPAMI.2004.75}

\bibitem[{Bron and Kerbosch(1973)}]{Bron1973}
Bron, C., Kerbosch, J.
\newblock Algorithm 457: finding all cliques of an undirected graph.
\newblock \emph{Communications of the ACM}, 16(9):575--577, 1973.
\newblock \href{https://doi.org/10.1145/362342.362367}{doi:10.1145/362342.362367}

\bibitem[{Prelog and Helmchen(1982)}]{Prelog1982}
Prelog, V., Helmchen, G.
\newblock Basic principles of the CIP-system and proposals for a revision.
\newblock \emph{Angewandte Chemie International Edition in English},
  21(8):567--583, 1982.
\newblock \href{https://doi.org/10.1002/anie.198205671}{doi:10.1002/anie.198205671}

\bibitem[{OpenSMILES(2024)}]{OpenSMILES2024}
OpenSMILES Working Group.
\newblock {OpenSMILES} specification.
\newblock \url{http://opensmiles.org/opensmiles-spec.pdf}, 2024.

\bibitem[{Bemis and Murcko(1996)}]{Bemis1996}
Bemis, G.W., Murcko, M.A.
\newblock The properties of known drugs. 1.~Molecular frameworks.
\newblock \emph{Journal of Medicinal Chemistry}, 39(15):2887--2893, 1996.
\newblock \href{https://doi.org/10.1021/jm9602928}{doi:10.1021/jm9602928}

\bibitem[{Steinbeck et~al.(2003)}]{Steinbeck2003}
Steinbeck, C., Han, Y., Kuhn, S., Horlacher, O., Luttmann, E., Willighagen, E.
\newblock The {Chemistry Development Kit (CDK)}: an open-source {Java} library
  for chemo- and bioinformatics.
\newblock \emph{Journal of Chemical Information and Computer Sciences},
  43(2):493--500, 2003.
\newblock \href{https://doi.org/10.1021/ci025584y}{doi:10.1021/ci025584y}

\bibitem[{Landrum(2006)}]{Landrum2006}
Landrum, G.
\newblock {RDKit}: open-source cheminformatics.
\newblock \url{https://www.rdkit.org/}, 2006.

\bibitem[{RDKit.js(2024)}]{RDKitJS2024}
RDKit.js Contributors.
\newblock {RDKit.js}: {RDKit} for {JavaScript} via {WebAssembly}.
\newblock \url{https://github.com/rdkit/rdkit-js}, 2024.

\end{thebibliography}

\end{document}
